\section{The appointed texts}
The Latin below follows the 1962 Missal, with spelling accents and ligatures
normalized and the printed abbreviated conclusions retained. Its preexisting
wording is witnessed in the public-domain books identified in the references.
The historical English prayers come from Cummiskey (1861), or from
Lasance--Walsh's revised missal for the three Maternity prayers.
Biblical English is Douay--Rheims--Challoner, in the Gutenberg 1581 witness.
These are study texts, not an approved English edition of the 1962 Missal.

The biblical extracts accompanying adapted chants are the cited Scripture,
not translations of the adapted Latin. The Introit repeats its antiphon
after the psalm verse and doxology in the ritual performance; its printed
abbreviation is retained here. The Maternity prayers apply only in the
Low or conventual branch described above.

\subsection{Introit}\label{proper-introit}

\properlocus{Daniel 3:31,29,35, adapted; Psalm 118:1 (119:1).}

\begin{latinproper}
Omnia, quae fecisti nobis, Domine, in vero iudicio fecisti, quia peccavimus tibi, et mandatis tuis non oboedivimus: sed da gloriam nomini tuo, et fac nobiscum secundum multitudinem misericordiae tuae. Ps. 118, 1 Beati immaculati in via: qui ambulant in lege Domini. V. Gloria Patri.
\end{latinproper}

\begin{englishwitness}{Biblical study English}

\textsuperscript{31} Wherefore, all that thou hast brought upon us, and every thing that thou hast done to us, thou hast done in true judgment:

\textsuperscript{29} For we have sinned, and committed iniquity, departing from thee: and we have trespassed in all things:

\textsuperscript{35} And take not away thy mercy from us, for the sake of Abraham, thy beloved, and Isaac, thy servant, and Israel, thy holy one:

\textsuperscript{1} Blessed are the undefiled in the way, who walk in the law of the Lord.

\end{englishwitness}

\properlocus{The antiphon rearranges and condenses Daniel; the cited verses above do not reproduce its complete verbal adaptation. Psalm 118:1 is the separate verse. Daniel 3:30,34,42--43 provides further verbal context.}

\subsection{Collect}\label{proper-collect}

\begin{latinproper}
Largire, quaesumus, Domine, fidelibus tuis indulgentiam placatus et pacem: ut pariter ab omnibus mundentur offensis, et secura tibi mente deserviant. Per Dominum.
\end{latinproper}

\begin{englishwitness}{Cummiskey, historical English}
Favourably grant, we beseech thee, O Lord, thy servants both pardon and peace; that, being cleansed from the guilt of all their offences, they may serve thee with secure minds. Thro'.
\end{englishwitness}

\subsection{Maternity Collect}\label{proper-maternity-collect}

\properlocus{Commemoration at Low or conventual Mass only.}

\begin{latinproper}
Deus, qui de beatae Mariae Virginis utero Verbum tuum, Angelo nuntiante, carnem suscipere voluisti: praesta supplicibus tuis; ut, qui vere eam Genetricem Dei credimus, eius apud te intercessionibus adiuvemur. Per eundem Dominum nostrum Iesum Christum Filium tuum: Qui tecum vivit et regnat in unitate.
\end{latinproper}

\begin{englishwitness}{Lasance--Walsh, historical English}
O God, Who wast pleased that at the angel's message Thy Word should take flesh in the womb of the blessed Virgin Mary, grant to Thy suppliants that, believing her to be truly the mother of God, we may be assisted by her intercessions with Thee. Through the same.
\end{englishwitness}

\subsection{Epistle}\label{proper-epistle}

\properlocus{Ephesians 5:15--21.}

\begin{latinproper}
Fratres: Videte quomodo caute ambuletis: non quasi insipientes, sed ut sapientes, redimentes tempus, quoniam dies mali sunt. Propterea nolite fieri imprudentes, sed intellegentes quae sit voluntas Dei. Et nolite inebriari vino, in quo est luxuria: sed implemini Spiritu Sancto, loquentes vobismetipsis in psalmis, et hymnis, et canticis spiritualibus, cantantes, et psallentes in cordibus vestris Domino: gratias agentes semper pro omnibus, in nomine Domini nostri Iesu Christi, Deo et Patri. Subiecti invicem in timore Christi.
\end{latinproper}

\begin{englishwitness}{Biblical study English}

\textsuperscript{15} See therefore, brethren, how you walk circumspectly: not as unwise,

\textsuperscript{16} But as wise: redeeming the time, because the days are evil.

\textsuperscript{17} Wherefore, become not unwise: but understanding what is the will of God.

\textsuperscript{18} And be not drunk with wine, wherein is luxury: but be ye filled with the Holy Spirit,

\textsuperscript{19} Speaking to yourselves in psalms and hymns and spiritual canticles, singing and making melody in your hearts to the Lord:

\textsuperscript{20} Giving thanks always for all things, in the name of our Lord Jesus Christ, to God and the Father:

\textsuperscript{21} Being subject one to another, in the fear of Christ.

\end{englishwitness}

\subsection{Gradual}\label{proper-gradual}

\properlocus{Psalm 144:15--16 (145:15--16).}

\begin{latinproper}
Oculi omnium in te sperant, Domine: et tu das illis escam in tempore opportuno. V. Aperis tu manum tuam: et imples omne animal benedictione.
\end{latinproper}

\begin{englishwitness}{Biblical study English}

\textsuperscript{15} The eyes of all hope in thee, O Lord: and thou givest them meat in due season.

\textsuperscript{16} Thou openest thy hand, and fillest with blessing every living creature.

\end{englishwitness}

\subsection{Alleluia}\label{proper-alleluia}

\properlocus{Psalm 107:2 (108:2), adapted.}

\begin{latinproper}
Alleluia, alleluia. V. Paratum cor meum, Deus, paratum cor meum: cantabo, et psallam tibi, gloria mea. Alleluia.
\end{latinproper}

\begin{englishwitness}{Biblical study English}

\textsuperscript{2} My heart is ready, O God, my heart is ready: I will sing, and will give praise, with my glory.

\end{englishwitness}

\properlocus{The chant has \emph{tibi, gloria mea}; the biblical English has ``with my glory.''}

\subsection{Gospel}\label{proper-gospel}

\properlocus{John 4:46b--53.}

\begin{latinproper}
In illo tempore: Erat quidam regulus, cuius filius infirmabatur Capharnaum. Hic cum audisset, quia Iesus adveniret a Iudaea in Galilaeam, abiit ad eum, et rogabat eum ut descenderet, et sanaret filium eius: incipiebat enim mori. Dixit ergo Iesus ad eum: Nisi signa et prodigia videritis, non creditis. Dicit ad eum regulus: Domine, descende priusquam moriatur filius meus. Dicit ei Iesus: Vade, filius tuus vivit. Credidit homo sermoni, quem dixit ei Iesus, et ibat. Iam autem eo descendente, servi occurrerunt ei, et nuntiaverunt dicentes, quia filius eius viveret. Interrogabat ergo horam ab eis, in qua melius habuerit. Et dixerunt ei: Quia heri hora septima reliquit eum febris. Cognovit ergo pater quia illa hora erat, in qua dixit ei Iesus: Filius tuus vivit: et credidit ipse, et domus eius tota.
\end{latinproper}

\begin{englishwitness}{Biblical study English}

\textsuperscript{46} He came again therefore into Cana of Galilee, where he made the water wine. And there was a certain ruler, whose son was sick at Capharnaum.

\textsuperscript{47} He having heard that Jesus was come from Judea into Galilee, went to him and prayed him to come down and heal his son: for he was at the point of death.

\textsuperscript{48} Jesus therefore said to him: Unless you see signs and wonders, you believe not.

\textsuperscript{49} The ruler saith to him: Lord, come down before that my son die.

\textsuperscript{50} Jesus saith to him: Go thy way. Thy son liveth. The man believed the word which Jesus said to him and went his way.

\textsuperscript{51} And as he was going down, his servants met him: and they brought word, saying, that his son lived.

\textsuperscript{52} He asked therefore of them the hour wherein he grew better. And they said to him: Yesterday at the seventh hour, the fever left him.

\textsuperscript{53} The father therefore knew that it was at the same hour that Jesus said to him: Thy son liveth. And himself believed, and his whole house.

\end{englishwitness}

\properlocus{The biblical English includes the Cana introduction in the first part of verse 46. The appointed Latin begins with the ruler in the latter part of that verse; verse 54 is outside the Gospel.}

\subsection{Offertory}\label{proper-offertory}

\properlocus{Psalm 136:1 (137:1), adapted.}

\begin{latinproper}
Super flumina Babylonis illic sedimus, et flevimus: dum recordaremur tui, Sion.
\end{latinproper}

\begin{englishwitness}{Biblical study English}

\textsuperscript{1} Upon the rivers of Babylon, there we sat and wept: when we remembered Sion:

\end{englishwitness}

\properlocus{The chant directly addresses Zion (\emph{tui}); the biblical verse names Zion as the object remembered.}

\subsection{Secret}\label{proper-secret}

\begin{latinproper}
Caelestem nobis praebeant haec mysteria, quaesumus, Domine, medicinam: et vitia nostri cordis expurgent. Per Dominum.
\end{latinproper}

\begin{englishwitness}{Cummiskey, historical English}
May these mysteries, O Lord, we beseech thee, procure us a heavenly remedy, and cleanse away the vices of our hearts. Thro'.
\end{englishwitness}

\subsection{Maternity Secret}\label{proper-maternity-secret}

\properlocus{Commemoration at Low or conventual Mass only.}

\begin{latinproper}
Tua, Domine, propitiatione, et beatae Mariae semper Virginis, Unigeniti tui Matris, intercessione, ad perpetuam atque praesentem haec oblatio nobis proficiat prosperitatem et pacem. Per eundem Dominum.
\end{latinproper}

\begin{englishwitness}{Lasance--Walsh, historical English}
Through Thy mercy, O Lord, and the intercession of blessed Mary, ever a virgin, the Mother of Thine only-begotten Son, may our oblation profit us for eternal and for present prosperity and peace. Through the same.
\end{englishwitness}

\subsection{Trinity Preface}\label{proper-preface}

\begin{latinproper}
Vere dignum et iustum est, aequum et salutare, nos tibi semper et ubique gratias agere: Domine, sancte Pater, omnipotens aeterne Deus: Qui cum unigenito Filio tuo, et Spiritu Sancto, unus es Deus, unus es Dominus: non in unius singularitate personae, sed in unius Trinitate substantiae. Quod enim de tua gloria, revelante te, credimus, hoc de Filio tuo, hoc de Spiritu Sancto, sine differentia discretionis sentimus. Ut in confessione verae sempiternaeque Deitatis, et in personis proprietas, et in essentia unitas, et in maiestate adoretur aequalitas. Quam laudant Angeli atque Archangeli, Cherubim quoque ac Seraphim: qui non cessant clamare cotidie, una voce dicentes:
\end{latinproper}

\begin{englishwitness}{Cummiskey, historical English}
It is truly meet and just, right and available to salvation, that we should always, and in all places, give thanks to thee, O holy Lord, Father Almighty, eternal God. Who together with thy only begotten Son and the Holy Ghost, art one God, and one Lord not in a singularity of one Person, but in a Trinity of one substance. For what we believe of thy glory, as thou hast revealed, the same we believe of thy Son, and of the Holy Ghost, without any difference or distinction. So that in the confession of the true and eternal Deity, we adore a distinction in the Persons, an unity in the essence, and an equality in the Majesty. Whom the angels and archangels, the cherubim also and seraphim praise: and cease not daily to cry out with one voice saying, Holy, \&c.
\end{englishwitness}

\properlocus{Cummiskey's common opening and Trinity continuation are joined according to his printed direction; his concluding ``Holy, \&c.'' points to the Sanctus of the Ordinary. The Latin Preface ends before that acclamation.}

\subsection{Communion}\label{proper-communion}

\properlocus{Psalm 118:49--50 (119:49--50), selected and adapted.}

\begin{latinproper}
Memento verbi tui servo tuo, Domine, in quo mihi spem dedisti: haec me consolata est in humilitate mea.
\end{latinproper}

\begin{englishwitness}{Biblical study English}

\textsuperscript{49} Be thou mindful of thy word to thy servant, in which thou hast given me hope.

\textsuperscript{50} This hath comforted me in my humiliation: because thy word hath enlivened me.

\end{englishwitness}

\properlocus{The chant ends at \emph{in humilitate mea}. The final life-giving clause of biblical verse 50 above is context, not part of the antiphon.}

\subsection{Postcommunion}\label{proper-postcommunion}

\begin{latinproper}
Ut sacris, Domine, reddamur digni muneribus: fac nos, quaesumus, tuis semper oboedire mandatis. Per Dominum nostrum.
\end{latinproper}

\begin{englishwitness}{Cummiskey, historical English}
That we may be worthy of thy sacred gifts, O Lord: grant, we beseech thee, we may always obey thy commandments. Thro'.
\end{englishwitness}

\subsection{Maternity Postcommunion}\label{proper-maternity-postcommunion}

\properlocus{Commemoration at Low or conventual Mass only.}

\begin{latinproper}
Haec nos communio, Domine, purget a crimine: et, intercedente beata Virgine Dei Genetrice Maria, caelestis remedii faciat esse consortes. Per eundem Dominum.
\end{latinproper}

\begin{englishwitness}{Lasance--Walsh, historical English}
May this communion, O Lord, purge away our guilt and, by the intercession of blessed Mary the Mother of God, make us companions of Him, Who is our heavenly healing. Through the same.
\end{englishwitness}

\properlocus{Lasance--Walsh render the heavenly remedy personally as Christ, ``Him, Who is our heavenly healing.'' This identification belongs to their English rendering.}
